\section{Conclusions}
\label{sec:conclusions}

This paper introduced a structured methodology for AI-assisted
conceptual aircraft design in which a large-language-model agent works
as an engineering collaborator inside a governed, version-controlled
design repository. The methodology rests on a strict separation of
roles --- deterministic, configured, regression-tested physics code as
the sole source of engineering numbers; the AI as the implementer and
propagator of design changes; humans as the source of requirements,
directives, review, and acceptance --- enforced by four transferable
mechanisms: a machine-readable project constitution, a one-way
analysis pipeline with typed hand-offs, pinned design-point regression
tests with full-pipeline continuous integration, and architecture
decision records with git-level provenance of AI contribution.

Alongside the workflow, the paper contributes an integrative parameter
model of electric VTOL UAV conceptual design: every design quantity is
classified as a requirement, a free parameter, a coupled state
variable, a derived quantity, or a constraint; the coupling map states
which factors influence which and locates the two circular
dependencies --- the mass closure and the stall-limited wing --- that
must be resolved simultaneously; the first-order dimensioning relations
that implement the map are given in full; and the constraint checks are
organised into three deliberately distinct failure exits, separating a
concept that does not close, a design that has drifted silently, and a
margin that is unfavourable but documented. This model is what makes
multidisciplinary change propagation a checkable operation rather than
an act of memory, and it is the component of a heritage design
capability that emerging drone organizations most conspicuously lack.
Its practical force is visible in the case study's own history, where
rotor diameter proved to be a mass parameter, battery internal
resistance a mass parameter, and maximum lift coefficient a geometry
parameter --- couplings that a sequential workflow reliably misses.

Applied to the conceptual design of a small electric tail-sitter VTOL
UAV, the methodology carried the program from mission requirements to
a converged, commercially frozen 2.52\,kg design point --- sixteen
analysis stages spanning sizing, aerodynamics, control authority,
vibration, thermal, CAD, component selection, and higher-fidelity
hand-off --- in three calendar weeks of part-time work, with thirteen
fully consistent tagged design snapshots and eighteen recorded
architecture decisions. Process evidence extracted from the git
history shows 64\% of design-phase commits co-authored by the AI, all
merged through human review, and documents which safety layer caught
each class of fault: automated pins caught silent drift, versioned
hand-off schemas caught cross-tool integration errors, and independent
expert review caught the conceptual errors no automated check could
see.

The answer to the research question is therefore affirmative but
conditional: AI can accelerate conceptual aircraft design by roughly
an order of magnitude in delivered scope per unit time \emph{when} it
is deployed at the workflow level under structural guardrails, rather
than as a source of engineering answers. The acceleration was used not
to cut corners but to make honesty cheap --- model refinements that
grew the design mass were adopted within days, and unfavorable margins
are carried openly as standing findings with named retirement paths.

Two consequences bear directly on practice. First, the framework
delivers not only a shorter development cycle but a higher quality of
work, because the dominant error mode of fast-moving programs --- a
parameter updated in one place and not in the places that depend on it
--- is removed from the space of possible outcomes rather than guarded
against by diligence. The benefit is therefore largest not for teams
with deep analysis heritage, whose main gain is speed, but for the far
larger population of integration-driven drone developers, who are
exposed precisely where the coupling map is dense and
counter-intuitive. Second, because the final pipeline stage regenerates
the design point, the margins table, the frozen hardware list, and
every standing finding on every run, the methodology doubles as an
instrument for formal design review: a package that cannot be stale by
construction, in which unmet requirements are structurally un-hideable,
every number is one step from the parameter and equation that produced
it, and a reviewer's alternative assumption can be re-converged during
the review rather than deferred as an action item.

Future work proceeds along the hand-offs the methodology already
exposes: closing the external vortex-lattice/blade-element and CFD
loops on the frozen geometry, PX4-based flight validation of the
hover/transition control architecture, folding measured component
properties (notably pack internal resistance) back into the closure
after procurement, and applying the framework to a second, dissimilar
airframe to test aircraft-independence. On the methodological side,
controlled multi-team studies are needed to separate the contributions
of the AI agent, the governance structure, and team expertise to the
observed acceleration.
